\section{Reception: Adamantius Among the Fathers}

\subsection{The ancient verdict}

The shape of Origen's reception was fixed by the fifth century: the greatest
of teachers, and the greatest of warnings. The nickname told the first half:
``He is said to have been also called Adamantius,'' Photius reports,
``because his arguments were linked together like chains of adamant''
(\work{Bibliotheca}, codex 118)---Jerome's ``man of adamant'' who out-wrote
Varro (\work{Letter} 33). Vincent of L\'erins, in 434, told both halves in
one chapter. Of the man: ``From his school came forth doctors, priests,
confessors, martyrs, without number''; ``what Christian did not reverence
him almost as a prophet; what philosopher as a master?'' Of the temptation:
who would not ``rather adopt that saying, That he would rather be wrong with
Origen, than be right with others''? And so, Vincent concludes, this Origen,
``such and so great as he was,'' proved for the Church ``not only a trial,
but a great trial'' (\work{Commonitorium} 17)---the same word he applied,
among the Latins, to Tertullian. The parallel is exact and was drawn in
antiquity: each language's supreme Christian intellect of the pre-Nicene
age, and neither on the calendar.

Origen is not a saint. He has no feast in any Roman liturgical book, no
recorded cult, no canonization, and---uniquely among figures of his rank in
Christian memory---a conciliar anathema attached to his name, in whatever
sense (see the preceding section). No commemoration of him, ancient or
modern, Eastern or Western, was located in the searches conducted for this
study; the medieval tomb tradition at Tyre attests fascination and respect,
not liturgical veneration (bounded negative results; see the appendix). Yet
he has never ceased to be treated as a Father in the working sense: the
patristic collections print him, the manuals cite him, and the Church's own
documents quote him---the \textit{Catechism} twice cites \work{On Prayer}
by name (CCC 2745, 2847). Citation, it should be stressed, is reception of
particular formulations judged sound, not rehabilitation of a corpus and
not a verdict on the man.

\subsection{Survival by translation, recovery by accident}

The sixth-century condemnations made Greek Origen a liability; copying
withered. What survived did so by four routes: the Greek of \work{Against
Celsus} and the \work{Philocalia} anthology, protected by their usefulness;
the Latin homilies and commentaries of Rufinus and Jerome; quotation in
friend and foe; and chance. Chance has been generous in the modern era. In
1941 wartime tunneling at Tura, near Cairo, uncovered papyri including
previously unknown Greek works (among them the \work{Dialogue with
Heraclides}, a transcript of Origen examining a bishop's theology before a
synod---the consulting theologian caught, for once, verbatim). In 2012
Marina Molin Pradel identified in a twelfth-century Munich manuscript
(Bayerische Staatsbibliothek, Cod.\ graec.\ 314) twenty-nine anonymous
Greek homilies on the Psalms as Origen's---the largest recovery of his
Greek preaching ever made, confirmed by Lorenzo Perrone's critical edition
(2015). The man burned by councils keeps returning by manuscript.

\subsection{The modern rehabilitation}

The twentieth century reversed the valence of the name. The patristic
revival around the series \textit{Sources chr\'etiennes}---founded in wartime
Lyon by Henri de Lubac and Jean Dani\'elou, and stocked from the first with
Origen's homilies---put the texts back in circulation. Dani\'elou's
\work{Orig\`ene} (1948) and above all de Lubac's \work{Histoire et esprit}
(1950) argued that Origen's ``spiritual'' exegesis was not arbitrary
allegory but a coherent theology of Scripture's unity in Christ, and made
him the fountainhead of the doctrine of the senses of Scripture that de
Lubac traced through the Middle Ages. Hans Urs von Balthasar's anthology
\work{Origenes: Geist und Feuer} (1938) had already called him the most
important theologian between Paul and Augustine. Henri Crouzel's lifetime
of studies, summed in his \work{Orig\`ene} (1985; English 1989), pressed
the historical case made first by Pamphilus: that Origen proposed where
later Origenism asserted, that he must be read as a man of the Church
exploring open questions, and that the condemned system is not his. Current
scholarship broadly accepts the distinction between Origen and
sixth-century Origenism while continuing to debate how far the difference
extends; the older English literature's ``Origen the heretic'' has largely
given way to ``Origen the first theologian.''

The papal reception followed the scholars. Benedict XVI's two general
audiences of 25 April and 2 May 2007 presented Origen as one of the
decisive figures of Christian thought: a ``true maestro'' whose pupils
could report that his life matched his doctrine; a man who fused theology
with exegesis so that thinking about God and interpreting Scripture became
one act, an ``irreversible turning point''; the pioneer of the threefold
reading of Scripture---letter, moral sense, spiritual mystery---and of the
praying study of the Bible that became \textit{lectio divina}; a teacher
whose insistence that knowledge of the Scriptures demands intimacy with
Christ in prayer ``concerns us all.'' Benedict noted the ardent youth who
longed to follow his father into martyrdom, the torture under Decius, and
the death from its effects; he did not treat the anathemas, and a general
audience neither could nor did reverse them. The juxtaposition is simply
the present state of the question: a pope commending, to the whole Church,
the spiritual doctrine of a man whose name stands in a conciliar list of
heretics. The Church has not resolved that tension; neither does this
study.
